\chapter{Améliorer la sélection par recherche locale et tabou}
\section{Un voisinage commun d'échange de stations}
À partir d'un ensemble $S$, on retire une station $u\in S\setminus\{r\}$ et
on ajoute un point $v\in V\setminus S$. Le voisin est défini par
\begin{equation}
S'=(S\setminus\{u\})\cup\{v\}.
\end{equation}
La cardinalité et la présence de la racine sont conservées par construction.
Le nombre de couples possibles vaut $(p-1)(n-p)$. Si $p=n$, aucun échange
n'existe et la recherche s'arrête avec un voisinage vide.

Pour chaque ensemble candidat, les affectations sont recalculées au plus proche,
le cycle est reconstruit par plus proche voisin puis amélioré par 2-opt.
L'évaluation porte sur le coût global $\alpha L+A$, et non sur le seul coût
d'accès. Cette reconstruction permet d'adapter le cycle à la nouvelle sélection,
mais son résultat reste heuristique. Elle ne calcule pas le meilleur cycle
possible pour chaque ensemble candidat.

\section{Recherche locale par échantillonnage}
La recherche locale part de la construction. À chaque itération, elle tire une
station retirable et un point à ajouter, évalue cet unique voisin, puis
l'accepte seulement s'il améliore strictement le coût courant. Un voisin rejeté
laisse la solution inchangée. La meilleure solution est donc aussi la solution
courante et son coût ne croît pas.

L'arrêt intervient lorsque le budget d'itérations ou de temps est atteint,
ou lorsque le voisinage est vide. Comme les voisins sont tirés au hasard et
peuvent être répétés d'une itération à l'autre, l'absence d'amélioration sur
les tirages observés ne prouve pas l'absence d'un meilleur voisin dans
l'ensemble complet. On parle ici de descente échantillonnée, et non de
recherche exhaustive jusqu'à certification d'un optimum local.

La graine initialise un générateur pseudo-aléatoire propre à la recherche.
Les listes de candidats sont triées avant le tirage pour stabiliser les choix.
Cette précaution évite de faire dépendre un résultat de l'ordre d'itération
d'un ensemble Python.

\section{Recherche tabou : accepter des détours contrôlés}
La recherche tabou distingue la solution courante de la meilleure solution
rencontrée. À chaque itération, elle échantillonne sans remise jusqu'à $k$
couples d'échange distincts dans le voisinage courant. Tous les candidats
évalués contribuent au compteur de travail, y compris ceux qui sont ensuite
écartés par la règle tabou.

Le mouvement admissible de plus faible coût est choisi dans cet échantillon.
Il peut dégrader le coût courant. Cette possibilité permet de quitter une
région où une descente stricte resterait bloquée. En parallèle, la meilleure
solution globale n'est remplacée que par une amélioration stricte : le résultat
retourné ne se dégrade donc pas avec les détours de la trajectoire.

Lorsqu'une station est retirée, sa réintroduction est interdite pendant la
durée tabou configurée. Une durée de un bloque bien son retour à l'itération
suivante. L'interdiction porte sur l'ajout du sommet retiré ; elle ne décrit pas
une mémoire de toutes les solutions visitées. Une règle d'aspiration autorise
un mouvement tabou s'il améliore la meilleure valeur globale.

S'il n'existe aucun mouvement admissible dans l'échantillon, la solution
courante reste inchangée et l'itération est comptée. Cela ne signifie pas que
tout le voisinage est interdit. Une nouvelle itération peut échantillonner
d'autres échanges ou bénéficier de l'expiration d'une interdiction.

\section{Déroulé de la recherche tabou}
\begin{enumerate}
\setlength{\itemsep}{1pt}\setlength{\parskip}{0pt}\setlength{\parsep}{0pt}
\item Initialiser la solution courante et la meilleure solution avec la construction.
\item Tirer jusqu'à $k$ échanges distincts et reconstruire leur solution complète.
\item Écarter les mouvements tabous, sauf en cas d'aspiration.
\item Adopter le meilleur mouvement admissible, même s'il dégrade la solution courante.
\item Mettre à jour les interdictions, la meilleure solution et les compteurs.
\item Répéter tant qu'un budget et un voisinage exploitable restent disponibles.
\end{enumerate}
Le nombre de mouvements effectués peut être inférieur au nombre d'itérations.
Ces deux nombres sont exportés séparément, de même que le nombre de candidats
évalués. Le budget est une limite demandée ; les compteurs décrivent le travail
effectivement réalisé.

\section{Budgets et limites de comparaison}
Dans la campagne principale, les deux méthodes disposent de 100 itérations.
La recherche locale évalue donc 100 voisins par exécution. La recherche tabou
évalue 15 voisins à chaque itération, soit 1\,500 dans les cas de cette campagne.
Les voisinages y sont assez grands pour fournir ces 15 couples distincts.
La comparaison porte ainsi sur deux réglages pratiques, avec des efforts différents.

Le budget temporel optionnel des recherches commence après la construction
initiale. Une évaluation déjà commencée peut dépasser légèrement la limite.
Les temps exportés pour l'exécution complète incluent, eux, les distances,
la construction et l'amélioration. Les horloges des traces sont relatives
à la phase d'amélioration ; elles n'ont donc pas exactement la même origine.

Des comparaisons à mêmes nombres d'évaluations et mêmes budgets de temps
permettraient de mieux isoler l'effet du tabou. La campagne temporelle
configurée dans le projet n'a pas été exécutée intégralement.
\bilan{Les deux recherches explorent la même famille d'échanges. La locale
privilégie les améliorations immédiates ; la tabou explore davantage et peut
franchir des dégradations temporaires. Les budgets expliquent une partie
du compromis observé entre coût final et temps.}
